\documentclass[11pt]{article}
\usepackage[T1]{fontenc}
\usepackage[utf8]{inputenc}
\usepackage{amsmath,amssymb,amsthm}
\usepackage[hmargin=0.9in,vmargin=0.85in]{geometry}
\usepackage{booktabs}
\usepackage{longtable}
\usepackage{url}
\usepackage[hidelinks]{hyperref}

\DeclareUrlCommand\dsix{\urlstyle{tt}}   % digraph6 strings, written \dsix+...+
\DeclareUrlCommand\cf{\urlstyle{tt}}     % file names in the project archive

\newtheorem{theorem}{Theorem}[section]
\newtheorem{lemma}[theorem]{Lemma}
\newtheorem{proposition}[theorem]{Proposition}
\newtheorem{corollary}[theorem]{Corollary}
\theoremstyle{definition}
\newtheorem{definition}[theorem]{Definition}
\newtheorem{remark}[theorem]{Remark}
\newtheorem{question}[theorem]{Question}

\newcommand{\R}{\mathbb{R}}
\newcommand{\Q}{\mathbb{Q}}
\newcommand{\one}{\mathbf{1}}
\newcommand{\eps}{\varepsilon}
\newcommand{\perpc}[1]{#1|_{\perp}}
\newcommand{\adj}{\operatorname{adj}}
\newcommand{\spec}{\operatorname{spec}}
\newcommand{\Rea}{\operatorname{Re}}
\newcommand{\lean}{\textsuperscript{\textsf{L}}}
\newcommand{\st}[1]{\textsc{#1}}

\title{The target-free clique conjecture for threshold-linear networks:\\
complete exact maps for $n\le 7$, infinite families,\\ and unbounded thresholds}
\author{}
\date{Draft, October 2026}

\begin{document}
\maketitle

\begin{abstract}
For combinatorial threshold-linear networks (CTLNs), Curto and Morrison asked whether the supports of
stable fixed points are exactly the target-free cliques of the underlying digraph. Geneson disproved this
with a six-node graph near $q=1$, where $q=\delta(1-\eps)/\eps>1$ is the legality parameter, and proved
that the conjecture holds once $q\ge n-2-(n-3)\eps/2$. We study the gap between these two results
with exact arithmetic throughout. First, a reduction theorem shows that at every legal parameter the
conjecture on $n$ nodes is equivalent to the statement that no non-complete digraph on $k\le n$ nodes
has a positive, linearly stable full-support fixed point. Second, using the rank-one structure of the
weight matrix, an exact $\eps\to0$ spectral criterion, Hurwitz determinants with a three-valued
treatment of boundary points, a cylindrical decomposition of the two-parameter legal region, and a
ray reduction to the legality curve, we obtain complete certified maps of stable non-clique full
supports: there are none for $n\le5$; at $n=6$ exactly three of the $1{,}540{,}944$ digraphs have one; at
$n=7$ exactly $124$ of the $882{,}033{,}440$ digraphs do, one of which is invisible in the $\eps\to0$ limit. The exact
thresholds are $q_6^*\approx1.1664$, a root of an explicit sextic, and $q_7^*=\rho\approx1.3247$, the
plastic number: for $n\le7$ the conjecture holds at every legal parameter with $q\ge\rho$ and fails for
every $q\in(1,\rho)$. Third, closed-twin blow-ups yield infinite families: counterexamples on every
$n\ge6$ nodes, finite-$\eps$-only counterexamples on every $n\ge7$ nodes, a proof that the critical
value $q_n^*$ tends to infinity, and a single legal parameter point at which stable non-clique supports
of every size $n\ge 6$ exist. Several of the necessary conditions used to prune the search are
machine-checked in Lean~4.
\end{abstract}

%======================================================================
\section{Introduction}\label{sec:intro}

\subsection{The conjecture}
A \emph{threshold-linear network} on $n$ neurons is the dynamical system
\begin{equation}\label{eq:tln}
\frac{dx_i}{dt}=-x_i+\Bigl[\sum_{j=1}^n W_{ij}x_j+\theta\Bigr]_+,\qquad i=1,\dots,n,
\end{equation}
with $[u]_+=\max(u,0)$ and a constant external drive $\theta>0$. A \emph{combinatorial} TLN (CTLN)
builds $W$ from a directed graph $G$ and two parameters $\eps,\delta>0$: weak inhibition $-1+\eps$
along edges and strong inhibition $-1-\delta$ along non-edges (Section~\ref{sec:prelim}). The CTLN is
\emph{legal} when $\delta>0$, $0<\eps<\delta/(\delta+1)$, which we write as $q:=\delta(1-\eps)/\eps>1$.
Fixed points of~\eqref{eq:tln} are indexed by their supports $\sigma\subseteq[n]$, and a central theme of
the CTLN programme is to read the fixed-point structure off the graph \cite{CurtoMorrison2023}.

A subset $\sigma$ is a \emph{clique} if every pair of distinct nodes of $\sigma$ is joined by edges in
both directions; a node $k\notin\sigma$ is a \emph{target} of $\sigma$ if every node of $\sigma$ sends an
edge to $k$; $\sigma$ is \emph{target-free} if it is a clique with no target. Open Question~10 of
Curto and Morrison \cite{CurtoMorrison2023} is the \emph{target-free clique conjecture}: \emph{the
supports of stable fixed points of a CTLN are exactly its target-free cliques.} One direction is a
classical graph rule: every target-free clique supports a stable fixed point at every legal parameter
(Lemma~\ref{lem:clique}). The content of the conjecture is the converse.

\subsection{Geneson's counterexample and bound}
Geneson \cite{Geneson2026} disproved the conjecture with an explicit graph on six nodes and
$\delta=\tfrac{29}{25}\eps$ for all sufficiently small $\eps$ (so $q\approx 29/25$), and proved
that it does hold whenever
\begin{equation}\label{eq:geneson-bound}
q\;\ge\; n-2-\tfrac{(n-3)\eps}{2}.
\end{equation}
This leaves a gap between $q$ slightly above $1$ and $q$ near $n-2$. The present paper maps that gap
exactly for $n\le7$ and studies its behaviour for large $n$.

\subsection{Our results}
Throughout, results are labelled \st{proved} (a proof is given here or in the project notes),
\st{certified} (an exact rational or algebraic computation recorded in a re-runnable certificate file),
or \st{Lean} (machine-checked in Lean~4 with Mathlib; marked \lean\ in statements). No floating-point
number enters any decision; decimals in the text are display approximations of exactly specified
algebraic numbers. Graphs on $n\le7$ nodes reported as counterexamples have been checked by hand (the two new
six-node graphs also by independent adversarial verification workflows, the $124$ seven-node graphs also by
four independently written exact verifiers); family members on $n\ge8$ nodes are certified instances
of proved theorems but have \emph{not} been hand-verified (Section~\ref{sec:verification}).

\paragraph{Reduction and exact machinery.} At every legal parameter, the conjecture on $n$ nodes is
equivalent to the absence of a non-complete digraph on $k\le n$ nodes whose full support is a positive,
linearly stable fixed point (Theorem~\ref{thm:reduction}). The full-support problem has a rank-one
structure: positivity depends on one parameter only, and the $\eps\to0$ limit is an exact spectral
condition on the adjacency matrix compressed to $\one^\perp$ (Lemmas~\ref{lem:rankone}
and~\ref{lem:limit}). Stability at some legal point is equivalent to stability on the one-dimensional
legality curve (Lemma~\ref{lem:rayR}). Together with five necessary conditions valid at every parameter
(domination, positivity, an edge bound, Stodola's condition and a twin obstruction, Section~\ref{sec:prunes})
these reduce the problem to a manageable set of exact one- and two-parameter certifications.

\paragraph{Complete maps for $n\le6$.} No digraph on at most five nodes is a counterexample at any legal
parameter. Among the $1{,}540{,}944$ digraphs on six nodes exactly three non-complete graphs have a
stable positive full support at some legal parameter: Geneson's graph, a partner obtained from it by a
bidirected in-twin switch, and one further 18-edge graph (Theorem~\ref{thm:n6}). The critical value is
$q_6^*\approx1.1664$, the root in $(561/481,7/6)$ of
$797q^6-1686q^5+2159q^4-3004q^3+2499q^2-2158q+1521$ (Theorem~\ref{thm:q6}).

\paragraph{The complete map at $n=7$.} Every one of the $882{,}033{,}440$ digraphs on seven nodes is either
excluded by a proved necessary condition or carries an exact two-parameter certificate. Exactly $124$
non-complete graphs have a stable positive full support at some legal parameter
(Theorem~\ref{thm:n7}). Of these, $123$ are visible in the $\eps\to0$ limit and one is not: its
compressed adjacency has spectral abscissa exactly $-1/2$ and it is stable only in a window of moderate
$\eps$. The natural monotonicity statement ``stable at finite $\eps$ implies stable in the limit'' is
false (Proposition~\ref{prop:Mfalse}), and an exact first-order coefficient $R^*$ decides when a stable
region bulges past its limit boundary (Theorem~\ref{thm:AS}). The threshold is $q_7^*=\rho$, the plastic
number (real root of $q^3=q+1$), and for $n\le7$ the conjecture holds at every legal parameter with
$q\ge\rho$ and fails for every $q\in(1,\rho)$ (Theorem~\ref{thm:q7}).

\paragraph{Infinite families and unbounded thresholds.} Closed-twin blow-ups preserve full-support
fixed-point data in closed form (Lemmas~\ref{lem:W} and~\ref{lem:X}). This yields: counterexamples on
every $n\ge6$ (Theorem~\ref{thm:famF}); finite-$\eps$-only counterexamples on every $n\ge7$
(Theorem~\ref{thm:famE}); families whose thresholds approach $1.84$, $3$ and $3.96$
(Theorem~\ref{thm:QQ2Y}); and, via an in-tree construction realising prescribed limit polynomials, a proof
that $q_n^*\to\infty$ (Theorem~\ref{thm:U}). Finally, at the single legal point
$(\eps,\delta)=(1/21,11/210)$ there are stable non-clique full supports on $n$ nodes for every
$n\ge6$, and on no $n\le5$ (Theorem~\ref{thm:A}).

\paragraph{Verification.} Every count and critical value is backed by a re-runnable certificate file
(Appendix~\ref{app:certs}); the pipelines and theorems were independently audited; all $n\le7$ counterexamples were
hand-verified, and the $124$ seven-node ones were re-decided by four independently written exact verifiers; and $87$ supporting theorems are machine-checked
(Section~\ref{sec:verification}).

\subsection{Related work}
The conjecture and the surrounding graph rules are from Curto and Morrison \cite{CurtoMorrison2023}.
Graphical domination, which we use as a prune in its verbatim form, is due to Curto \cite{Curto2025}.
Geneson \cite{Geneson2026} gave the first counterexample and the bound~\eqref{eq:geneson-bound}; he does
not pose or treat the gap question.
An unrefereed preprint, ``Six Is Minimal'' \cite{SixIsMinimal2026}, obtained independently and by
different exact methods that no digraph on at most five nodes has a stable non-clique support, used
Geneson's graph as the six-node witness, computed the full-support region of that graph, and gave a
family on $5r+6$ nodes. Their conventions agree with ours up to a transpose of the non-edge matrix; their
exclusion for $n\le5$ agrees with our $n=5$ map; their region for Geneson's graph gives
$1+u^*=1.16635502785\ldots$, the same sextic root as our $q_6^*$; and their family is a clique-class
blow-up of the same seed with $q\to1$, an instance of Lemma~\ref{lem:W}. Their work does not classify
$n=6$ or treat $n\ge7$, finite-$\eps$-only examples, the reduction theorem or growing thresholds; it does
give an explicit rational parameter schedule for its family and global nondegeneracy certificates, which
prompted our nondegeneracy check. We regard the two works as independent and consistent.
Our comparison is based on reading their repository; we did not run their code.

%======================================================================
\section{Preliminaries}\label{sec:prelim}

\subsection{Conventions}
A digraph $G$ on $[n]$ is encoded by its adjacency matrix $A$, with
\[
A_{ij}=1\iff\text{there is an edge } i\to j\qquad(\text{row } i \text{ lists the nodes that } i \text{ sends to}).
\]
The CTLN weight matrix is
\begin{equation}\label{eq:W}
W_{ij}=\begin{cases}0,& i=j,\\ -1+\eps,& i\neq j,\ A_{ji}=1\quad(j\to i:\ \text{weak inhibition}),\\
-1-\delta,& i\neq j,\ A_{ji}=0\quad(j\not\to i:\ \text{strong inhibition}).\end{cases}
\end{equation}
Note the transpose: $W_{ij}$ is the influence of $j$ on $i$ and therefore reads $A_{ji}$. For a support
$\sigma$ of size $m$ with induced adjacency $A_\sigma$, with $J=\one\one^{T}$ the all-ones matrix,
\begin{equation}\label{eq:closed}
I-W_\sigma\;=\;-\delta I+(1+\delta)J-(\eps+\delta)A_\sigma^{T}.
\end{equation}
Every entry is a polynomial in $(\eps,\delta)$ with integer coefficients. Identity~\eqref{eq:closed},
including the transpose, is machine-checked (\texttt{closed\_form}\lean).

\begin{definition}\label{def:sfp}
A nonempty $\sigma\subseteq[n]$ \emph{supports a stable fixed point} (SFP) if
\begin{enumerate}
\item[(1)] $(I-W_\sigma)x=\theta\one$ has a solution with $x>0$ entrywise;
\item[(2)] $\theta+\sum_{j\in\sigma}W_{kj}x_j<0$ for every $k\notin\sigma$;
\item[(3)] every eigenvalue of $-I+W_\sigma$ has negative real part.
\end{enumerate}
\end{definition}
If (1)--(3) hold, the zero-extension of $x$ is an equilibrium of~\eqref{eq:tln} near which the vector field is
affine with Jacobian spectrum $\spec(-I+W_\sigma)\cup\{-1\}$, so it is linearly, hence locally
asymptotically, stable. ``Stable'' always means linearly stable in this strict sense: points where
$I-W_\sigma$ is singular or a Hurwitz determinant vanishes are never counted as stable. When
$\sigma=[n]$, condition~(2) is vacuous and ``target-free clique'' reduces to ``complete bidirected''.

Condition (3) is decided by the Hurwitz determinants of the exact characteristic polynomial: for a real
$f(z)=c_0z^m+\dots+c_m$ with $c_0>0$, $\Delta_k(f)=\det[c_{2i-j}]_{i,j=1}^{k}$ (with $c_l=0$ outside
$0\le l\le m$), and $f$ has all roots in $\Rea z<0$ iff $\Delta_1,\dots,\Delta_m>0$ (Routh--Hurwitz;
see Gantmacher, \emph{The Theory of Matrices}, Vol.~II, Ch.~XV).

\subsection{Reparametrisation and legality}
Put
\[
t=\frac{1-\eps}{\eps},\qquad \beta=1+\frac{\delta}{\eps},\qquad \gamma=1-\frac1\beta,\qquad
q=\frac{\delta(1-\eps)}{\eps}=\frac{(\beta-1)t}{t+1}.
\]
The map $(\eps,\delta)\mapsto(\beta,t)$ is a homeomorphism of $(0,1)\times(0,\infty)$ onto
$(1,\infty)\times(0,\infty)$ with inverse $\eps=1/(1+t)$, $\delta=(\beta-1)/(1+t)$, and
\begin{equation}\label{eq:legal}
\text{legal}\iff q>1\iff t(\beta-2)>1
\end{equation}
(\texttt{legal\_iff}\lean). In particular every legal point has $\beta>2$ and $\gamma\in(1/2,1)$, and
along a ray $\delta=(\beta-1)\eps$, $\eps\to0$, one has $t\to\infty$ and $q\to\beta-1$. The drive
$\theta>0$ only rescales fixed points: each of (1)--(3), and every boundary event, has the same truth
value at every $\theta>0$ (Lemma~G of the reduction note\footnote{\cf{notes/lemma_F_reduction.md}, \S1.}).

Let $N:=J-I-A_\sigma^T$ be the \emph{receiver-first non-edge matrix}: $N_{ab}=1$ iff $a\ne b$ and
$\sigma_b\not\to\sigma_a$. Put $C:=I+\beta N$ and $K:=C+tJ$.

\subsection{Rank-one structure}
\begin{lemma}[Rank-one structure; proved]\label{lem:rankone}
Let $\eps>0$. Then
\begin{enumerate}
\item[(i)] $I-W_\sigma=\eps K$ (\texttt{IminusW\_eq}\lean);
\item[(ii)] $a:=\adj(K)\one=\adj(C)\one$ for every value of $t$; each $a_i\in\mathbb Z[\beta]$ has degree $\le m-1$;
\item[(iii)] $\det K=d+ts$ with $d=\det C$ and $s=\one^Ta$;
\item[(iv)] $\det(\mu I+K)=\phi(\mu)+t\,\omega(\mu)$ with $\phi(\mu)=\det(\mu I+C)$ and
$\omega(\mu)=\one^T\adj(\mu I+C)\one$, where $\omega$ has degree exactly $m-1$ in $\mu$ and leading coefficient $m$;
\item[(v)] if $D:=\det K\ne0$, the solution of (1) is $x=(\theta/\eps)\,a/D$, and for $k\notin\sigma$ the
drive in (2) equals $\theta\, o_k/D$ with $o_k=d-\beta\sum_b N_{kb}a_b$, where $N_{kb}=1$ iff $\sigma_b\not\to k$.
\end{enumerate}
\end{lemma}
\begin{proof}[Sketch]
(i) is an entrywise check using $\eps\beta=\eps+\delta$ and $\eps t=1-\eps$. (ii)--(iv) follow from the
adjugate form of the matrix determinant lemma, $\det(M+\tau\one\one^T)=\det M+\tau\,\one^T\adj(M)\one$,
valid over any commutative ring, and Cramer's rule in column form: in $K$ with column $i$ replaced by
$\one$, subtracting $t$ times that column from the others gives $C$ with column $i$ replaced by $\one$.
(v) is Cramer's rule together with $W_{kb}=-\eps(t+\beta N_{kb})$. No step divides by $\det C$.
Full proofs: \cf{notes/lemmas.md}, \S1.
\end{proof}

\begin{lemma}[Positivity is $t$-free; proved, \lean]\label{lem:A2}
If $K$ is positive stable (equivalently, (3) holds) then $\det K>0$. Hence, at every point with $\eps>0$,
\[
\sigma\text{ is an SFP}\iff a_i(\beta)>0\ \forall i,\quad o_k(\beta)<0\ \forall k\notin\sigma,\quad
\det(\mu I+K)\ \text{Hurwitz}.
\]
\end{lemma}
\begin{proof}[Sketch]
$\det K$ is the product of the eigenvalues; real ones are positive and complex ones pair to
$|\lambda|^2>0$. The Lean proof (\texttt{det\_pos\_of\_posStable}, \texttt{ctln\_pos\_iff\_adjugate\_pos})
argues instead that $x\mapsto\det(xI+K)$ is monic with no root in $[0,\infty)$.
\end{proof}
The Lean development covers $\det K>0$ and the positivity clause ($x>0\iff a>0$ under stability); the off-support
clause $o_k<0$ is not formalised.
Thus in the $(\beta,t)$ quarter-plane the SFP set is a product of a set of $\beta$-intervals (positivity
and the off-support condition) with $(0,\infty)$, intersected with the stability set and the legal
region; only stability depends on $t$.

\subsection{The \texorpdfstring{$\eps\to0$}{eps -> 0} limit}
Let $U\in\R^{m\times(m-1)}$ have orthonormal columns spanning $\one^\perp$ and write
$\perpc{X}:=U^TXU$. Let $\alpha:=\max\Rea\spec(\perpc{A_\sigma})$ be the spectral abscissa of the
compressed adjacency; it depends on the graph only. Exactly, $\spec\perpc{A}$ can be computed over $\Q$
from any rational basis of $\one^\perp$.

\begin{lemma}[Limit criterion; proved]\label{lem:limit}
Let $\beta>1$ and $\gamma=1-1/\beta$.
\begin{enumerate}
\item[(a)] $\omega(\mu;\beta)=m\det\bigl((\mu+1)I_{m-1}+\beta\,\perpc{N}\bigr)$, and as $t\to\infty$ one eigenvalue
of $K(\beta,t)$ is $mt+O(1)$ while the other $m-1$ converge to $\spec\bigl((1-\beta)I-\beta\perpc{A}^T\bigr)$.
\item[(b)] If $\alpha<-\gamma$ there is $T$ with $K(\beta,t)$ positive stable for all $t>T$; if
$\alpha>-\gamma$, then $K(\beta,t)$ is unstable for all large $t$; $\alpha=-\gamma$ is undecided by the limit.
The limit-stable set of legal $\beta$ is the interval $(2,\beta_c)$, $\beta_c=1/(1+\alpha)$, nonempty iff
$\alpha<-1/2$; always $\alpha\ge-1$, with equality iff $G[\sigma]$ is complete bidirected. The threshold $T$ is
uniform on compact subsets of $(2,\beta_c)$.
\item[(c)] $a(\beta)$ is one-signed $\iff$ $A^T+\gamma I$ is nonsingular and $w:=(A^T+\gamma I)^{-1}\one>0$; then $a$ is a
nonzero multiple of $w$, and condition (1) holds at $(\beta,t)$ for some, equivalently every, $t>0$ iff $w>0$. For $k\notin\sigma$,
$o_ks<0$ iff $\sum_{b:\sigma_b\to k}w_b<1$.
\item[(d)] If $\alpha<-\gamma$ at a legal $\beta$, then $G[\sigma]$ has more than $m(m-1)/2$ edges.
\end{enumerate}
\end{lemma}
\begin{proof}[Sketch]
(a) Conjugating by the orthogonal matrix $[\one/\sqrt m\,|\,U]$ turns $tJ$ into $tm\,e_1e_1^T$, and the
determinant is affine in the $(1,1)$ entry with slope the complementary minor; root continuity (Rouch\'e,
or an elementary continuity lemma) gives the eigenvalue statement. (b) follows from (a) and
$\perpc{J}=0$, since $\perpc{N}=-I-\perpc{A}^T$. (c) uses $Ca=d\one$ and $C=-\beta(A^T+\gamma I)+\beta J$, so
$(A^T+\gamma I)a=\frac{\beta s-d}{\beta}\one$, together with $\det(A^T+\gamma I)=(-1/\beta)^m(d-\beta s)$. (d) is
a trace computation: $\operatorname{tr}\perpc{A}=-|E|/m$ and the maximum real part dominates the mean.
Full proofs: \cf{notes/lemmas.md}, \S4, and \cf{notes/lemma_I_spectral.md}, \S\S2--6.
\end{proof}

Combining (b) and (c): if $w(\gamma)>0$ and $\alpha<-\gamma$ at some rational $\beta>2$, then the full support
is an SFP at the legal points $(\eps,\delta)=(1/(1+t),(\beta-1)/(1+t))$ for all $t>T(\beta)$. We call such a
graph \emph{limit-stable}. Lemma~\ref{lem:limit} is the shared dependency of every ``for all $k$''
statement in Section~\ref{sec:families}; it is not formalised in Lean. Parts (b) and (c) are Lemmas~I.3 and~I.4 of
the project notes; they were refereed when the $n=6$ map was audited and were not re-derived by the referees of the
family theorems.

\subsection{Reduction to full supports}
\begin{lemma}[Cliques; proved]\label{lem:clique}
Let $0<\eps<1$, $\delta>0$, $\theta>0$ and let $\sigma$ be a clique of size $m$. Then $x=c\one$ with
$c=\theta/(\eps+m(1-\eps))$ solves (1), (3) holds, and the drive at $k\notin\sigma$ is
$\theta[\eps-r(\eps+\delta)]/(\eps+m(1-\eps))$ where $r$ is the number of $j\in\sigma$ with $j\not\to k$.
Hence a clique is an SFP iff it is target-free.
\end{lemma}

For a legal point $p$ and $k\ge2$ let $\mathrm{FS}(k,p)$ be the statement ``every $k$-node digraph whose full
support satisfies (1) and (3) at $p$ is complete bidirected''.

\begin{theorem}[Reduction to full supports; proved]\label{thm:reduction}
Fix a legal $p=(\eps,\delta)$ and $n\ge1$. The following are equivalent:
\begin{enumerate}
\item[(i)] the target-free clique conjecture holds for every $n$-node CTLN at $p$;
\item[(ii)] for every $n$-node $G$, every SFP support of $G$ at $p$ is a clique;
\item[(iii)] $\mathrm{FS}(k,p)$ holds for every $k\in\{2,\dots,n\}$.
\end{enumerate}
\end{theorem}
\begin{proof}[Sketch]
(i)$\Leftrightarrow$(ii) by Lemma~\ref{lem:clique}. (iii)$\Rightarrow$(ii): conditions (1) and (3) for
$\sigma$ in $G$ involve only $W_\sigma$, which is the weight matrix of the induced subgraph $G[\sigma]$; so a
non-clique SFP $\sigma$ would give a non-complete $|\sigma|$-node graph satisfying (1) and (3) on its full
support. (ii)$\Rightarrow$(iii) by padding: given a non-complete $H$ on $k\le n$ nodes with (1) and (3) on
its full support, add $n-k$ nodes receiving no edges from $H$. At such a node the drive is
$\theta-\sum_jc_{kj}x_j$ with all $c_{kj}=1+\delta$; subtracting row $i$ of the fixed-point equation
shows it is at most $-\delta x_i<0$, so (2) holds and $V(H)$ is a non-clique SFP of the padded graph.
Condition (2) enters only to exclude targets of cliques. Full proof:
\cf{notes/lemma_F_reduction.md}, \S6 (there called Theorem~F).
\end{proof}

\begin{corollary}\label{cor:monotone}
Let $\mathrm{Bad}(n)$ be the set of legal points at which the conjecture fails for some $n$-node graph, and $S_H$
the set of legal points at which the full support of $H$ is an SFP. Then
$\mathrm{Bad}(n)=\bigcup_{k\le n}\bigcup_{H}S_H$ over non-complete $k$-node $H$, and
$\mathrm{Bad}(n)\subseteq\mathrm{Bad}(n+1)$.
\end{corollary}

\begin{lemma}[Openness; proved]\label{lem:open}
Each $S_H$ is open in the $(\eps,\delta)$ plane. Consequently every nonempty $S_H$ contains rational points,
and at every point of $S_H$ all Hurwitz determinants and coefficients of the characteristic polynomial are
strictly positive: a counterexample point is never an ``undecided'' point.
\end{lemma}

We write $q_n^*$ for the supremum of $q$ over legal points at which some non-complete $n$-node digraph has a
stable positive full-support fixed point. By Corollary~\ref{cor:monotone}, the conjecture for all graphs
on at most $n$ nodes holds at every legal point with $q\ge\max_{k\le n}q_k^*$ and fails at points with $q$
arbitrarily close to that maximum. For $k\le4$ there is no counterexample at any legal parameter: $k=2$ by
a direct eigenvalue computation (for a one-way edge $\det K=1-t(\beta-2)$, negative exactly at legal points; a non-adjacent pair is excluded by Lemma~\ref{lem:twin}), $k=3$ by Theorem~4.1 of
\cite{Geneson2026}, $k=4$ is stated as known in \cite{Geneson2026} (for nondegenerate CTLNs, citing work of Curto, Geneson and
Morrison whose exact hypotheses we have not checked), and all three cases are re-certified by our two-parameter engine on all
digraphs with at most four nodes; our claims for $k\le4$ rest on that certification.\footnote{\cf{certificates/region2d_n2to4_all.json}; \cf{notes/lemma_F_reduction.md}, \S\S5, 8.}

%======================================================================
\section{The exact machinery}\label{sec:machinery}

\subsection{Boundary sets and three-valued decisions}
At a single sampled parameter a vanishing Hurwitz determinant says nothing about the neighbourhood, so
every pointwise procedure in this work returns one of three values: \emph{stable}, \emph{not stable}, or
\emph{undecided}. An undecided sample is never counted as unstable and never as a counterexample. Regions
are decided instead on open cells of a decomposition whose walls are the zero sets of a proved boundary set.

\begin{lemma}[Imaginary-axis roots; proved]\label{lem:b3}
If $f$ is real of degree $m\ge2$ and $f(iw)=0$ for some real $w\neq0$, then $\Delta_{m-1}(f)=0$.
\end{lemma}
\begin{proof}[Sketch]
With $P(x)=x^mf(1/x)$, the factor $1+w^2x^2$ divides $P=(1+w^2x^2)R$; for $V(x)=xR(-x)\neq0$ the product $PV$ is odd,
and its even coefficients are exactly the entries of $H_{m-1}v$ where $H_{m-1}$ is the matrix whose determinant is
$\Delta_{m-1}$. This is the special case of Orlando's formula that is needed; it is proved self-contained in
\cf{notes/lemmas.md}, \S2.
\end{proof}

\begin{lemma}[Boundary set at fixed $t$; proved]\label{lem:B}
Fix $t>0$ and let $p_t(\mu;\beta)=\det(\mu I+K)$. Let $P_t$ be the family
$a_1,\dots,a_m$, $D_t=d+ts$, $o_k$ ($k\notin\sigma$), $\Delta_{m-1}(p_t)$, all in $\Q[\beta]$. If some member is
identically zero, $\sigma$ is an SFP at no legal $\beta$. Otherwise the SFP verdict is constant on each component of the
legal $\beta$-range minus the real roots of $P_t$, and false at those roots. The same holds with $\beta$ fixed and $t$
varying, and in the limit $t=\infty$ with $\omega$, $s$, $\Delta_{m-2}(\omega)$ in place of $p_t$, $D_t$,
$\Delta_{m-1}(p_t)$.
\end{lemma}
\begin{proof}[Sketch]
On a component, every member of $P_t$ has constant sign; by Lemma~\ref{lem:b3} no root of $p_t$ meets the imaginary
axis; by continuity of roots the number of roots in the right half-plane is locally constant. \cf{notes/lemmas.md}, \S2.
\end{proof}

The two-parameter version (Lemma~B$'$ and Corollaries~B$''$, B$'''$ of the project
notes\footnote{\cf{notes/lemma_Bprime_2d.md}.}) shows that the verdict is invariant on connected components of
the complement of the curves $\{a_i=0\}$, $\{D=0\}$, $\{\Delta_{m-1}=0\}$ and the legality boundary
$L(\beta,t)=t(\beta-2)-1=0$ in the open legal region, and that over each open $\beta$-interval free of the roots of a
projection set (leading coefficients, discriminants and pairwise resultants of these polynomials in $t$) the sectors
of a cylindrical algebraic decomposition are sign-invariant cells. Each full-dimensional cell is decided at one exact
rational sample; lower-dimensional cells are boundary points and are never SFP.

\subsection{The engines}
\paragraph{One-parameter engine.} At fixed $t$ (or $t=\infty$), the polynomials of Lemma~\ref{lem:B} are obtained by
exact integer evaluation and interpolation with provably sufficient degree bounds, their real roots are isolated by
continued fractions into irreducible minimal polynomials with rational isolating intervals, and each cell is decided
at an exact sample by the Hurwitz routine. Every isolation is replayed by Sturm counts and every sample is
cross-checked against an independent pointwise verifier. The engine agrees with an independent symbolic path on
all $7{,}194$ (graph, support, $\eps$) cases at $n=4$ and on $900$ at $n=5$.\footnote{\cf{certificates/validate_fastregion_n4.json}, \cf{certificates/validate_fastregion_n5.json}.}

\paragraph{Two-parameter engine.} A cylindrical algebraic decomposition of the legal $(\beta,t)$ plane on the proved
boundary set, restricted to the positivity cells, with a Descartes-bisection replay of every root count that is
independent of the isolator. It reproduces the theorem for $k\le4$ on all digraphs with at most four nodes and agrees
with an earlier implementation on all $706$ five-node survivors of a weaker filter.\footnote{\cf{certificates/region2d_n2to4_all.json}, \cf{certificates/region2d_v2_validation.json}.}

\paragraph{Policy.} A cell whose Hurwitz test is undecided at three samples is \emph{escalated}, not decided; there were $0$
escalations at $n=5,6,7$. By Lemma~\ref{lem:open} undecided points cannot hide counterexamples.

\subsection{Ray reduction to the legality curve}
\begin{lemma}[Ray identity; proved, \lean]\label{lem:rayF}
$K(\lambda\beta,\lambda t)=(1-\lambda)I+\lambda K(\beta,t)$. Hence if $K(\beta,t)$ is positive stable, so is
$K(\lambda\beta,\lambda t)$ for every $\lambda\in(0,1]$ (\texttt{affine\_identity}, \texttt{posStable\_ray}\lean).
\end{lemma}

\begin{lemma}[Ray reduction; proved, forward direction \lean]\label{lem:rayR}
The full support is stable at some legal $(\eps,\delta)$ if and only if $K(2+1/t,\,t)$ is positive stable for some
$t>0$. Moreover, if the full support is an SFP at a legal $(\beta_0,t_0)$, the corresponding curve point has
$\beta'<\beta_0$.
\end{lemma}
\begin{proof}
($\Leftarrow$) Positive stability is open, so $K(2+1/t+\eta,t)$ is stable for small $\eta>0$, and that point is legal.
($\Rightarrow$) Let $(\beta,t)$ be legal and stable and $f(\lambda)=\lambda t(\lambda\beta-2)-1$. Then $f(2/\beta)=-1<0<f(1)$, so
$f$ has a root $\lambda^*\in(2/\beta,1)$, explicitly $\lambda^*=(t+\sqrt{t^2+\beta t})/(\beta t)$; the point
$(\lambda^*\beta,\lambda^*t)$ lies on the curve $t(\beta-2)=1$ and is stable by Lemma~\ref{lem:rayF}.
The forward direction is \texttt{lemmaR\_forward\_ctln}\lean; the converse needs continuity of eigenvalues and is not
formalised.
\end{proof}

Since positivity is $t$-free (Lemma~\ref{lem:A2}), a graph that is never stable on the curve has no SFP anywhere in
the legal plane. Along the curve, $tK(2+1/t,t)=L(t):=tI+(2t+1)N+t^2J$ is an integer polynomial matrix; the coefficient
of $\mu^{m-k}$ in $\det(\mu I+L(t))$ has $t$-degree at most $k+1$, and by Lemma~\ref{lem:B} with $t$ as the parameter
the verdict changes only at positive roots of $\det L(t)$ and of $\Delta_{m-1}(t)$. One exact sample per root-free
interval decides the curve (the \emph{curve test}). A fast path discards a graph when no $t>0$ makes every coefficient
positive (Stodola's condition, below). The curve test takes about 35--45~ms per seven-node graph, and on all $9{,}608$
five-node digraphs returns $95$ stable and $9{,}513$ never, each of the $95$ with an exact legal witness.\footnote{\cf{n7/validate_curve_n5.json}; \cf{notes/lemma_M_refuted.md}, \S3.}

\subsection{Necessary conditions valid at every parameter}\label{sec:prunes}
Every prune below is a theorem valid at \emph{every} legal parameter, and all five are machine-checked. The twin
obstruction is not a separate stage of the $n=7$ chain.

\begin{lemma}[Domination; proved, \lean]\label{lem:D}
Node $k$ \emph{graphically dominates} $j\ne k$ if $j\to k$, $k\not\to j$, and every $i\notin\{j,k\}$ with $i\to j$ also has
$i\to k$ \cite[Def.~1.1]{Curto2025}. If $k$ dominates $j$, then every fixed point of~\eqref{eq:tln} with uniform drive has
$x_j=0$, for every $\eps>0$ and $\delta\ge0$; in particular the full support has no positive fixed point. The same holds
for domination inside an induced subgraph $G[\sigma]$ (\texttt{lemmaD}, \texttt{lemmaD\_full\_support}, \texttt{lemmaD\_induced}\lean).
\end{lemma}
\begin{proof}[Sketch]
For $x\ge0$ one has $y_j(x)-y_k(x)\le(1-\eps)x_j-(1+\delta)x_k$, where $y_l$ is the input to node $l$. At a fixed
point with $x_j>0$ this forces $y_k>0$ and then $0\ge\eps x_j+\delta x_k>0$. Neither $\eps<1$ nor legality is used, but
the uniform drive is: with $\theta_j>\theta_k$ the conclusion fails. \cf{notes/lemmas.md}, \S5.
\end{proof}

\begin{lemma}[Edge bound; proved, \lean]\label{lem:e2}
Let $B$, $N_n$, $O$ be the numbers of bidirected, non-adjacent and one-way unordered pairs of an $m$-node graph. The
sum of the $2\times2$ principal minors of $K$ is
\[
e_2(K)=B(1+2t)+N_n(1-\beta)(1+\beta+2t)+O\bigl(1-t(\beta-2)\bigr).
\]
If $K$ is positive stable at a legal point then $e_2(K)>0$, hence $B>N_n$ and the graph has at least $m(m-1)/2+1$ edges
(\texttt{e2\_closed\_form}, \texttt{edge\_bound\_ctln}\lean).
\end{lemma}
At a legal point the non-adjacent and one-way terms are negative, which gives $B>N_n$. In Lean the spectral step
(positive stability forces $e_2>0$) is proved from Mathlib's description of characteristic-polynomial coefficients as
sums of principal minors and the Gauss--Lucas theorem. For $n=7$ the bound reads ``at least $22$ edges''.

\begin{lemma}[Stodola; proved, \lean]\label{lem:stodola}
A monic real polynomial with all roots in $\Rea z<0$ has all coefficients positive. Hence if $K$ is positive stable,
every coefficient of $\det(\mu I+K)$ is positive (\texttt{stodola}, \texttt{coeff\_pos\_of\_posStable}, \texttt{ctln\_coeff\_pos}\lean).
\end{lemma}

\begin{lemma}[Twin obstruction; proved, \lean]\label{lem:twin}
If $u,v\in\sigma$ are non-adjacent and have the same in-neighbours in $\sigma\setminus\{u,v\}$, then $e_u-e_v$ is a left
eigenvector of $W_\sigma$ with eigenvalue $1+\delta$, so $-I+W_\sigma$ has the eigenvalue $\delta>0$ and $\sigma$ is never
stable, for every $\eps$ and every $\delta>0$ (\texttt{twin\_not\_hurwitz}, \texttt{twin\_not\_hurwitz\_induced}\lean). For
bidirected in-twins the eigenvalue is $-\eps$.
\end{lemma}

Together with Lemma~\ref{lem:A2} (a positive stable full support needs a legal $\beta$ with all $a_i(\beta)>0$) and the
density bound Lemma~\ref{lem:limit}(d) for the limit, these are the only prunes used. An early placeholder for domination
was replaced by the verbatim definition of \cite{Curto2025} before any sweep used it.

%======================================================================
\section{Complete maps for \texorpdfstring{$n\le6$}{n <= 6}}\label{sec:n6}

\subsection{Validation gates}
Two published results were reproduced before any new computation was trusted.

\emph{Geneson's counterexample.} The paper \cite{Geneson2026} gives its graph as a receiver-first non-edge matrix; our
adjacency matrix is its complement and transpose, and the weight convention is the same. The printed integer
certificate $Cv=179298949\cdot\one$ reproduces exactly, our fixed point is proportional to the paper's $v$, and the
full support is a stable non-clique fixed point at $\eps=1/100$, $\delta=29/2500$ and at $\eps=1/1000$,
$\delta=29\eps/25$, while it is a positive but unstable fixed point at $\eps=1/10,1/20,1/50$ on the same ray.\footnote{\cf{certificates/stage0_geneson_2607.21396.json}; a permanent test in the project's test suite.}

\emph{The published $n=5$ classification} (the package accompanying Curto's $n=5$ data, which stores the transpose of
our adjacency matrix). At its default parameters $\eps=1/4$, $\delta=1/2$, every recorded fixed-point support of every
one of the $9{,}608$ five-node digraphs matches our exact verifier: $24{,}442$ supports, $0$ disagreements, $0$ boundary
cases (the package records no stability data; our verifier finds $14{,}488$ of the supports stable). Under the wrong (untransposed) convention a sampled version of this check fails on
$27$ of $30$ graphs.\footnote{\cf{certificates/stage1_n5_diff.json}.}

Our enumerations reproduce the number of digraphs up to isomorphism (OEIS A000273): $9{,}608$ for $n=5$,
$1{,}540{,}944$ for $n=6$ and $882{,}033{,}440$ for $n=7$.

\subsection{\texorpdfstring{$n\le5$}{n <= 5}}
\begin{theorem}[$n\le5$; certified]\label{thm:n5}
No non-complete digraph on at most five nodes has a stable positive full-support fixed point at any legal parameter.
Consequently, by Theorem~\ref{thm:reduction}, the target-free clique conjecture holds for every CTLN on at most five
nodes at every legal parameter.
\end{theorem}
Of the $9{,}608$ five-node digraphs, $390$ admit a legal $\beta$ with all $a_i(\beta)>0$; the two-parameter
decomposition certifies $389$ of them as never stable and only $K_5$ as stable, with $0$ escalations. (An earlier,
weaker common-sign filter left a superset of $706$, which was also certified in full: $705$ never, $K_5$
stable.)\footnote{\cf{certificates/region2d_n5_shard0of1_tally.json}, \cf{certificates/n5_summary.json}.}
This agrees with the independent result of \cite{SixIsMinimal2026}.

\subsection{\texorpdfstring{$n=6$}{n = 6}}
\begin{theorem}[$n=6$; certified, hand-verified]\label{thm:n6}
Among the $1{,}540{,}944$ digraphs on six nodes, exactly three non-complete graphs have a stable positive full-support
fixed point at some legal parameter (Table~\ref{tab:n6}). Every other six-node graph is certified never stable and
positive at every legal $(\eps,\delta)$.
\end{theorem}

\begin{table}[ht]
\centering
\footnotesize
\begin{tabular}{@{}llccl@{}}
\toprule
graph & adjacency rows & edges & limit-stable $q$ & relation\\
\midrule
$G_1$ (index $1{,}000{,}255$) & \texttt{011000 000011 000111 110011 110101 101110} & 19 & $(1,q_6^*)$ & Geneson\\
$G_2$ (index $1{,}424{,}629$) & \texttt{011000 000101 000111 100011 110101 111110} & 19 & $(1,q_6^*)$ & switch of $G_1$\\
$H_3$ (index $1{,}387{,}528$) & \texttt{011000 000011 000111 110010 100101 111110} & 18 & $(1,q_c')$ & new\\
\bottomrule
\end{tabular}
\caption{The three six-node counterexamples, in canonical labelling (row $i$ lists the $j$ with $i\to j$); $G_2$ is a
bidirected in-twin switch of $G_1$. $q_c'\approx1.0830$ is the root in $(1,13/12)$ of
$q^5+2q^4+q^3-2q^2-2q-1$.}\label{tab:n6}
\end{table}

\paragraph{The chain.} In the limit sweep over all $1{,}540{,}944$ graphs, $1{,}512{,}594$ have no legal $\beta$ with all
$a_i>0$ and therefore (Lemma~\ref{lem:A2}) no stable positive full support at any parameter; $28{,}350$ strict survivors
remain, all among the $204{,}378$ domination-free graphs; exactly four are limit-stable ($K_6$, $G_1$, $G_2$, $H_3$);
$0$ escalations. The two-parameter decomposition was run on a superset of $34{,}071$ survivors (an earlier weaker
filter): $34{,}067$ never stable anywhere, $4$ stable somewhere, $0$ escalations, $0$ errors, every cell cross-checked
pointwise and every root count replayed. Exact $\beta$-slices at $\eps=1/1000$ and $1/100$ show $K_6$ and the three
graphs; at $\eps=1/10$ only $K_6$ and Geneson's class; at $\eps=1/4$ and $1/2$ only $K_6$.\footnote{\cf{certificates/n6_summary.json}; \cf{certificates/region2d_n6_shard{0,1,2}of3*}; \cf{certificates/region_n6_eps*}.}

\paragraph{The regions.} For $G_1$ and $G_2$ the stable positive region is
$\{2<\beta<\beta_c,\ t>\max(1/(\beta-2),T(\beta))\}$ with $T(\beta)<1/(\beta-2)$ on an initial interval $(2,\beta_1)$,
$\beta_1$ a degree-$14$ algebraic number in $(17/8,15/7)$: these graphs are stable at \emph{every} legal $t$ when $\delta/\eps$
is close enough to $1$, so counterexamples exist with $q$ arbitrarily close to $1$. For $H_3$ the two thresholds
coincide: the full support is an SFP iff $2<\beta<\beta_1'$, at every legal $t$, where $\beta_1'\approx2.08302$ is the root
in $(152/73,25/12)$ of $\beta^5-3\beta^4+3\beta^3-3\beta^2+2\beta-1$, a positivity boundary ($a_5=0$). A witness is
$\eps=1/101$, $\delta=27/2626$, $q=1350/1313$.\footnote{\cf{certificates/COUNTEREXAMPLE_CANDIDATE_n6_idx1387528.json}, \cf{certificates/COUNTEREXAMPLE_CANDIDATE_n6_idx1424629.json}.}

\paragraph{Twins.} $G_2$ is obtained from $G_1$ by moving one out-edge between a pair of bidirected in-twins.

\begin{proposition}[Bidirected in-twin switching; proved]\label{prop:J4}
Let $u\leftrightarrow v$ be bidirected and have the same in-neighbours outside $\{u,v\}$, and let $G'$ be obtained from $G$
by replacing an edge $u\to w$ (with $v\not\to w$) by $v\to w$. Then for every support, every $\eps,\delta>0$ and every
$\theta$, $G$ and $G'$ have the same fixed points (same vectors) and the same Jacobian characteristic polynomials.
\end{proposition}
More generally, two graphs have identical full-support data for all parameters iff after relabelling they are cospectral
and have the same in-walk counts $(A^T)^k\one$ for all $k$; the $34{,}071$ six-node survivors form $21{,}800$ such classes.
Bidirected in-twin switching preserves all supports, whereas other members of a class can differ in their proper
supports.\footnote{\cf{notes/lemma_J_equivalence.md}; \cf{certificates/n6_survivor_equivalence_classes.json}.}
So modulo twin switching there are exactly two six-node counterexample graphs.

\subsection{The threshold at \texorpdfstring{$n\le6$}{n <= 6}}
In the limit, $\beta_c(G_1)=-1/\Rea\nu_1$, where $\nu_1$ is a root of the quartic factor $6z^4+11z^3+16z^2+10z+3$ of the
characteristic polynomial of the compressed non-edge matrix (the other factor is $z$, from a pair of twins). Real parts of
roots of a quartic are roots of a polynomial of degree at most $\binom42=6$; here $\beta_c$ is a root of the irreducible sextic
$797b^6-6468b^5+22544b^4-44440b^3+53280b^2-38016b+13824$, and the substitution $b=q+1$ gives the sextic in $q$ below.

\begin{theorem}[Threshold for $n\le6$; certified]\label{thm:q6}
Let $q_6^*$ be the root in $(561/481,\,7/6)$ of
\[
797q^6-1686q^5+2159q^4-3004q^3+2499q^2-2158q+1521\qquad(q_6^*\approx1.1664).
\]
For every graph on at most six nodes, the target-free clique conjecture holds at every legal parameter with
$q\ge q_6^*$; for every $q\in(1,q_6^*)$ there is a legal parameter with that value of $q$ at which it fails for a six-node
graph. The supremum is approached by $G_1$ and $G_2$ as $t\to\infty$ and is not attained.
\end{theorem}
\begin{proof}[Sketch]
By Theorem~\ref{thm:n5}, Theorem~\ref{thm:n6} and Corollary~\ref{cor:monotone}, $\mathrm{Bad}(6)=S_{G_1}\cup S_{G_2}\cup S_{H_3}$. On every
finite-$t$ cell $q=(\beta-1)t/(t+1)<\beta-1$, and the exact suprema over the certified cells are computed by the
threshold tool from the two-parameter certificates; the limit interval of $G_1$ starts at $q=1$, so the union covers
$(1,q_6^*)$.\footnote{\cf{certificates/qstar/qstar_n6.json}, tool \cf{n7/qstar.py}.}
\end{proof}
Geneson's bound~\eqref{eq:geneson-bound} at $n=6$ reads $q\ge4-3\eps/2$; Theorem~\ref{thm:q6} shows the true value is
$q_6^*\approx1.1664$. The same sextic root, $1.16635502785\ldots$, was obtained independently in \cite{SixIsMinimal2026}.

%======================================================================
\section{The complete map at \texorpdfstring{$n=7$}{n = 7}}\label{sec:n7}

\subsection{The accounting chain}
\begin{theorem}[$n=7$; certified, hand-verified]\label{thm:n7}
Among the $882{,}033{,}440$ digraphs on seven nodes, exactly $124$ non-complete graphs have a stable positive full-support
fixed point at some legal $(\eps,\delta)$. Each carries an exact two-parameter region certificate. Of the $124$, exactly
$123$ are limit-stable; the remaining graph, \dsix+&FNFL~[udo?+, is stable only at finite $\eps$.
\end{theorem}

Every digraph is either removed by a necessary condition valid at every legal parameter, or certified by the
two-parameter engine (Table~\ref{tab:chain}). Each stage was validated before use against the $n=6$ map (the positivity
stages reproduce the $28{,}350$ six-node strict survivors exactly, with zero difference in both directions; the curve
test applied to those survivors returns exactly the four stable six-node graphs).

\begin{table}[ht]
\centering
\small
\begin{tabular}{@{}cp{0.52\textwidth}rr@{}}
\toprule
stage & condition (where proved) & removed & remaining\\
\midrule
0 & enumeration (OEIS A000273) & & $882{,}033{,}440$\\
1 & a dominated node (Lemma~\ref{lem:D}) & $764{,}800{,}518$ & $117{,}232{,}922$\\
2 & no $\beta>2$ with all $a_i(\beta)>0$, exact Descartes bounds (Lemma~\ref{lem:A2}) & $107{,}600{,}789$ & $9{,}632{,}133$\\
3 & undecided lines of stage 2 resolved by exact root isolation & $723{,}010$ & $8{,}909{,}123$\\
4 & at most $21$ edges (Lemma~\ref{lem:e2}) & $5{,}135{,}775$ & $3{,}773{,}348$\\
5 & never stable on the legality curve (Lemma~\ref{lem:rayR}) & $3{,}773{,}183$ & $165$\\
6 & stable somewhere, never stable and positive together (CAD) & $40$ & $125$\\
\bottomrule
\end{tabular}
\caption{Accounting chain at $n=7$. The $125$ survivors are $K_7$ and the $124$ counterexamples; $0$ escalations at every
stage.}\label{tab:chain}
\end{table}

Of the $9{,}632{,}133$ lines kept at stage 2, $8{,}845{,}537$ were certified positive and $786{,}596$ were undecided by the
fast exact filter; stage 3 resolved the latter exactly. The single run performing stages 4 and 5 over the $8{,}909{,}123$ positivity survivors took $16.4$ CPU-hours ($6.6$ ms per line).\footnote{\cf{certificates/n7_finite_summary.json} (field \texttt{chain}, \texttt{accounting\_ok = true}); \cf{certificates/n7_curve_TOTAL.json}; \cf{certificates/n7_region2d/}; runbook \cf{n7/REPRODUCE_n7_finite.md}.}
Independently, the limit classification was run first: a spectral prefilter over all $882{,}033{,}440$ digraphs left $638{,}042$
survivors ($1{,}889$ strict), whose exact limit classification gave the $123$ limit-stable non-clique graphs and
$K_7$.\footnote{\cf{certificates/COUNTEREXAMPLE_CANDIDATES_n7_limit.jsonl}, \cf{certificates/n7_limit_tally.json}. For $15$ graphs the first
stored pointwise witness in the limit file is an illegal point ($q\le1$); use \cf{certificates/n7_candidates_legal_witnesses.json}.}
The finite-$\eps$ map contains all of them, as Lemma~\ref{lem:limit} requires.

\subsection{The 124 graphs}
\begin{itemize}
\item Edge counts range from $23$ to $31$ of $42$ possible (counts $7,7,18,10,6,22,34,17,3$ for $23,\dots,31$ edges); all
satisfy the edge bound ($\ge22$) and all are domination-free.
\item They form $91$ classes of identical full-support data (Section~\ref{sec:n6}): $71$ singletons, $16$ pairs, two classes of
$4$, one of $5$ and one of $8$; $64$ graphs carry a bidirected in-twin pair.
\item Of the $123$ limit-stable graphs, the limit critical values $q_c$ (right ends of the limit-stable $q$-intervals) have $69$ distinct minimal polynomials;
the most frequent is $q^5+2q^4+q^3-2q^2-2q-1$ ($17$ graphs, the same quintic as for $H_3$), and the largest $q_c$ is the plastic
number, the real root of $q^3-q-1$ ($6$ graphs). For $6$ of the $123$ the positivity set starts strictly above $\beta=2$.
\item $21$ of the $123$ contain a six-node counterexample as an induced subgraph; none contains $K_6$.
\end{itemize}
Every graph comes with a legal exact witness point.\footnote{\cf{certificates/COUNTEREXAMPLE_CANDIDATES_n7_finite.jsonl}; \cf{certificates/n7_candidates_legal_witnesses.json}.}

\begin{corollary}[certified]\label{cor:alpha7}
Every seven-node graph whose full support is an SFP at some legal parameter has $\alpha\le-1/2$, with equality only for
\dsix+&FNFL~[udo?+.
\end{corollary}

\subsection{A counterexample invisible in the limit}
The graph $E_0=$ \dsix+&FNFL~[udo?+ ($28$ edges) has adjacency rows \texttt{0011110} \texttt{0011100} \texttt{1101111}
\texttt{1110111} \texttt{0011011} \texttt{0100101} \texttt{1100000}. Its compressed adjacency satisfies
\[
7\chi_{\perpc{A}}(\lambda)=(\lambda+1)(\lambda^2+\lambda+1)(7\lambda^3+14\lambda^2+14\lambda+5),
\]
so $\alpha=-1/2$ exactly and, by Lemma~\ref{lem:limit}(b), the graph is limit-stable at no legal $\beta$. Nevertheless its full
support is a stable positive fixed point at $\eps=1/49$, $\delta=33/1568$ ($q=99/98$) and at $\eps=1/36$, $\delta=17/588$. Its
stable region is bounded in $t$ at every $\beta$: on its single positivity cell $2<\beta<\beta_1'$ ($\beta_1'\approx2.08302$, the quintic root of Table~\ref{tab:n6},
isolated here in $(27/13,25/12)$), at the sample $\beta=53/26$, the stable range is
$26<t<t_1$ with $t_1$ an algebraic number in $(53,54)$, and the full support is unstable for larger $t$. The supremum of $q$
over its region is an algebraic number of degree $25$ in $(685003/657781,\,2170535/2084278)$, about $1.0414$, approached at the
corner where $\Delta_6=0$ meets the positivity boundary, not as $t\to\infty$.\footnote{\cf{certificates/COUNTEREXAMPLE_CANDIDATE_n7_finite_FNFL.json}; \cf{certificates/qstar/qstar_fnfl.json} (the exact bracket is recorded there and in \cf{qstar_n7_finite.json}; the file's aggregate block is stale and reports only a weaker sample bound).}
Structurally, $E_0$ is $H_3$ with one node doubled into a closed-twin pair, and the factor $\lambda^2+\lambda+1$ is a \emph{moving}
eigenvalue pair (not an eigenvalue of $A$) that sits exactly on the boundary line; see Theorem~\ref{thm:famE}.

\subsection{The limit does not determine the finite map}
Write $c=1+t/\beta>1$. A one-line computation gives the identity
\begin{equation}\label{eq:E}
K=\beta\,(cJ-A^T-\gamma I),
\end{equation}
so $K$ is positive stable iff $\nu_A(c):=\min\Rea\spec(cJ-A)>\gamma$, and legality is $\gamma>c-\sqrt{c^2-c}$, which is the
stability boundary of a single one-way edge. As $c\to\infty$, $\nu_A(c)\to-\alpha$. The natural hope is that
$\nu_A$ is nondecreasing, i.e.\ that stability at a legal finite $t$ implies stability for all larger $t$ at the same $\beta$
(``Lemma~M'').

\begin{proposition}[Lemma M is false; certified]\label{prop:Mfalse}
(a) For $A$ with rows \texttt{011 001 100} and $\beta=22/5$, $K$ is positive stable at the legal point $t=11/25$ and not at
$t=1$ or $t=100$; here $\alpha=-2/3>-17/22=-\gamma$.
(b) For the domination-free seven-node graph \dsix+&FK@dyUfXp?+ (one of the $123$) at $\beta=10561/5000$, the full support is
an SFP at $t=10$ and positive but unstable at $t=1000$, and the limit polynomial is not Hurwitz.
\end{proposition}
Example (a) is not a CTLN counterexample (it fails positivity), but (b) shows the failure persists with positivity. In the
latter, the right end of the stable $\beta$-interval at $t=10,30,100,1000,10^4$ lies strictly above the limit value
$\beta_c$.\footnote{\cf{notes/lemma_M_refuted.md}, \S2; scripts in \cf{notes/prover_scripts_M/}.}
Lemma~\ref{lem:A2} shows that every such failure under positivity is an imaginary-pair crossing: on the positivity set
$s>0$, so $D=d+ts$ increases in $t$ and a real eigenvalue crosses $0$ only upward.

\subsection{The bulge criterion}
Let $\chi(z)=\det(zI+A)$, $\psi(z)=\one^T\adj(zI+A)\one$ and $\hat s=\psi/\chi=\one^T(zI+A)^{-1}\one$. Then
$\det(zI-(cJ-A))=\chi(z)-c\psi(z)$; roots of $\gcd(\chi,\psi)$ are eigenvalues for every $c$ (\emph{fixed}), the others are
\emph{moving}.

\begin{theorem}[First-order asymptotics; proved]\label{thm:AS}
Assume $\nu_\infty:=-\alpha\in(1/2,1)$ and that every root $z_0$ of $\psi$ with $\Rea z_0=\nu_\infty$ is simple and not a root of
$\chi$. Let $R^*=\min\Rea\bigl(1/\hat s\,'(z_0)\bigr)$ over these roots. Then $\nu_A(c)=\nu_\infty+R^*/c+O(c^{-2})$, and the
boundary of the stable region near $(\beta_c,\infty)$ is $\beta=\beta_c+\beta_c^3R^*/t+O(t^{-2})$.
\begin{enumerate}
\item[(a)] If $R^*>0$, then for every $\beta$ slightly above $\beta_c$ the stable $t$-set is nonempty and bounded above; if
moreover $a(\beta_c)>0$, these points are SFPs.
\item[(b)] If $R^*<0$, the stable $t$-set is upward closed for large $t$ near $\beta_c$.
\end{enumerate}
\end{theorem}
\begin{proof}[Sketch]
With $e=1/c$ the moving eigenvalues solve $\psi_1(z)-e\chi_1(z)=0$ ($\chi=h\chi_1$, $\psi=h\psi_1$, $h=\gcd$); the implicit function
theorem gives an analytic branch $z(e)=z_0+e\,\chi_1(z_0)/\psi_1'(z_0)+O(e^2)$, and $\chi_1(z_0)/\psi_1'(z_0)=1/\hat s\,'(z_0)$.
\cf{notes/lemma_M_refuted.md}, \S4.
\end{proof}

\begin{proposition}[Bulge rule at $n=7$; certified]\label{prop:bulge}
For all $124$ graphs, $R^*$ was computed exactly ($51$ positive, $73$ negative, none zero; every critical root simple and not
shared with $\chi$). A graph has a stable cell bounded above in $t$ if and only if $R^*>0$ and its fixed point is positive at
$\beta_c$; this holds for exactly $40$ graphs, and every bounded cell starts exactly at $\beta_c$. The $11$ graphs with $R^*>0$
that lose positivity before $\beta_c$ do not bulge as fixed points. For $E_0$, $R^*=1/4$ and $\beta_c=2$.\footnote{\cf{certificates/n7_bulge_Rstar_exact.json}; \cf{notes/bulge_exact/}.}
\end{proposition}

At $n=6$ no stable cell is bounded in $t$ (Geneson's class has $R^*<0$, and the limit interval of $H_3$ ends at a positivity
root), so the $\eps\to0$ classification describes the six-node regions faithfully, while it understates the regions of $40$ of
the seven-node graphs and misses $E_0$ entirely.

\subsection{The threshold at \texorpdfstring{$n\le7$}{n <= 7}}
\begin{theorem}[Threshold for $n\le7$; certified, hand-verified]\label{thm:q7}
Let $\rho\approx1.3247$ be the plastic number, the real root of $q^3-q-1$ (isolated in $(45/34,4/3)$). Then $q_7^*=\rho$, and
the supremum is not attained. Consequently, for every graph on at most seven nodes the target-free clique conjecture holds at
every legal parameter with $q\ge\rho$, and for every $q\in(1,\rho)$ there is a legal parameter with that value of $q$ at which it
fails for a seven-node graph.
\end{theorem}
\begin{proof}[Sketch]
On every finite cell $q<\beta-1$. On the $258$ stable cells unbounded in $t$ the supremum is $\beta_{\mathrm{hi}}-1$ with
$\beta_{\mathrm{hi}}$ a limit endpoint, and the largest is $1+\rho$ (six graphs, whose limit interval ends at the positivity root
$\beta^3-3\beta^2+2\beta-1=0$, i.e.\ $q^3-q-1=0$, with the stability boundary further right). On the $40$ bounded cells the outer
isolating bound gives $\beta_{\mathrm{hi}}\le16/7$, so $q<9/7<\rho$. Four of the six plastic-number graphs have limit intervals
starting at $q=1$, so the union of regions covers $(1,\rho)$. Combine with Theorem~\ref{thm:q6} and
Corollary~\ref{cor:monotone}. An independent run of the exact threshold tool over all $298$ stable cells returns matching
upper and lower bounds with no undecided cell.\footnote{\cf{certificates/qstar/qstar_n7_finite.json}; \cf{certificates/qstar/qstar_n7limit.json}.}
\end{proof}

Geneson's bound gives $q\ge5-2\eps$ at $n=7$; the true threshold is the plastic number. The clique-class family of
Theorem~\ref{thm:famF} alone reaches only about $1.2747$ at $n=7$, and other graphs approach the supremum.

%======================================================================
\section{Infinite families}\label{sec:families}

The theorems in this section are proved for every member of a family; tables and individual instances are certified
only. Instances on $n\ge8$ nodes are certified by the pointwise exact verifier but have \emph{not} been hand-verified, and
are reported as candidates in that sense. The proofs use the base graphs only as explicit matrices. Theorems~\ref{thm:famF},
\ref{thm:famE}, \ref{thm:QQ2Y} and~\ref{thm:U} rest on Lemma~\ref{lem:limit}(b),(c) (Lemmas~I.3 and~I.4 of the notes:
transfer from the limit to finite $t$, and $t$-free positivity); Theorem~\ref{thm:A} uses only part~(c).

\subsection{The non-edge picture and blow-ups}
Let $D:=J-I-A$ be the sender-first non-edge matrix ($D=N^T$ for the full support) and
$\psi_D(z):=\one^T\adj(zI-D)\one=m\,\chi_{\perpc{D}}(z)$. Since $\perpc{A}=-I-\perpc{D}$, Lemma~\ref{lem:limit} reads:
\emph{the full support is limit-stable at $\beta$ iff every root of $\psi_D$ has real part $>-1/\beta$}; and the fixed point is
positive iff $u:=(D^T+\beta^{-1}I)^{-1}\one>0$ and $\one^Tu>1$, in which case $w=u/(\one^Tu-1)$.

\begin{lemma}[Closed-twin blow-up; proved]\label{lem:W}
Let $H$ have $m$ nodes and let $\mu\in\mathbb Z_{>0}^m$, $M=\operatorname{diag}(\mu)$. Let $H^\mu$ replace node $a$ by a bidirected
clique of $\mu_a$ nodes, with every node of class $a$ sending to every node of class $b\ne a$ iff $a\to b$ in $H$; $n=\sum\mu_a$.
Then
\[
n\,\chi_{\perpc{D'}}(z)=z^{n-m}\,\mu^T\adj(zI-DM)\one,
\]
and for every $\gamma$ with $A'^T+\gamma I$ nonsingular the positivity vector of $H^\mu$ is constant on classes, $w'=Ey$ with
$\bigl((A^T+I)M+(\gamma-1)I\bigr)y=\one$, where $E$ is the class-indicator matrix.
\end{lemma}
\begin{proof}[Sketch]
$D'=EDE^T$ and $E^TE=M$; push-through $(zI-EDE^T)^{-1}E=E(zI-DM)^{-1}$ and Sylvester's identity
$\det(zI_n-EDE^T)=z^{n-m}\det(zI_m-DM)$.\footnote{\cf{notes/research_wf3/theory/notes_theory.md}, \S2.}
\end{proof}

For a base $G$, in-set $S$ and out-set $T$, let $G[S,T;k]$ add a bidirected clique $V$ of $k$ nodes, with $i\to v$ for all $v\in V$
iff $S_i=1$ and $v\to i$ iff $T_i=1$; put $b=\one-S$, $c=\one-T$.

\begin{lemma}[Clique-class extension; proved]\label{lem:X}
(X1) $n\,\chi_{\perpc{D'}}=z^{k-1}\Phi_k$ with $\Phi_k(z)=z\,\psi_D(z)+k\,\chi_L(z)$, $L:=D-b\one^T-\one c^T$.
(X2) If $T$ is all of $G$, then for every $k\ge1$ and rational $\beta>1$ the full support of $G[S,T;k]$ is positive iff
$w_G(\beta)>0$ and $\sigma_S(\beta):=\sum_{i\in S}w_{G,i}(\beta)>1$; this does not depend on $k$ or $t$.
(X3) As $k\to\infty$ the moving spectrum converges to $\spec(L)$. With several sink classes of sizes $k_j$ and in-sets $S_j$,
the same holds with $L$ replaced by $D-\bar b\one^T$, $\bar b$ the size-weighted mean of the $b_j$, and positivity holds iff
$w_G>0$ and $\sigma_{S_j}>1$ for every $j$.
\end{lemma}

\subsection{Counterexamples on every \texorpdfstring{$n\ge6$}{n >= 6}}
Let $G_1$ be Geneson's graph in the paper's labelling (rows \texttt{011011 101011 110101 111000 011000 000110}), $S=(1,1,1,1,1,0)$,
and $F_k:=G_1[S,\text{all};k]$, $n=6+k$.

\begin{theorem}[Clique-class family; proved, refereed]\label{thm:famF}
For every $k\ge0$, $\alpha(\perpc{F_k})<-1/2$, and $F_k$ is positive at every $\beta\in(2,b_p)$, where $b_p-1=\rho$ is the
plastic number. Hence for each $k$ and each rational $\beta\in(2,\min(b_p,\beta_c(k)))$ the full support of $F_k$ is a stable
positive non-clique fixed point at $(\eps,\delta)=(1/(1+t),(\beta-1)/(1+t))$ for all large $t$. In particular the target-free
clique conjecture fails for every $n\ge6$.
\end{theorem}
\begin{proof}[Sketch]
By Lemma~\ref{lem:X}, $\Phi_k=z^2R_k$ with $R_k=(6+k)z^4+(11+k)z^3+(16+k)z^2+10z+3$. The shifted polynomial
$16R_k(-\tfrac12-x)=16(k+6)x^4+16(k+1)x^3+8(2k+17)x^2+12(k+1)x+(3k+16)$ has Hurwitz determinants
$16(k+1)$, $64(k+1)(k+16)$, $8192(k+1)^2$ and $8192(k+1)^2(3k+16)$, all positive for $k>-1$. Positivity: the components of
$\adj(\beta A^T+(\beta-1)I)\one$ and $\sigma_S-1$ are explicit polynomials; exact root counting on $[2,93/40]$ shows a single
sign change, at the root $b_p$ of $b^3-3b^2+2b-1$.\footnote{\cf{notes/research_wf3/theory/family_F_certificate.json} (script \cf{family_F_proof.py}); witnesses for $k=1,\dots,6,8,10$ ($n=7,\dots,16$) at $(\eps,\delta)=(1/21,11/210)$ or $(1/21,11/175)$.}
\end{proof}
The limit supremum of $q$ along the family is $\min(b_p,\beta_c(k))-1$: about $1.1664$ ($k=0$), $1.2747$ ($k=1$), the plastic
number $1.3247$ for $k=2,\dots,8$ ($n=8,\dots,14$), then decreasing to $1$ (since $\chi_L=z^4(z^2+z+1)$ has roots on
$\Rea z=-1/2$). The same construction proves $9$ families on the three six-node graphs and $617$ families on $117$ of the
$124$ seven-node graphs, each for every $k\ge1$, by exact real-root isolation in $k$; an independent referee re-proved all $626$
for every real $k\ge1$.\footnote{\cf{notes/research_wf3/theory/families_catalogue_n6.json}, \cf{families_catalogue_n7.json}.}

\subsection{Finite-\texorpdfstring{$\eps$}{eps}-only counterexamples on every \texorpdfstring{$n\ge7$}{n >= 7}}
\begin{lemma}[Boundary mechanism; proved]\label{lem:Bmech}
Suppose (i) $\psi$ has a conjugate pair of simple roots $z_0,\bar z_0$ with $\Rea z_0=1/2$ and $\chi(z_0)\ne0$; (ii) every other
root of $\psi$ has real part $>1/2$; (iii) $R:=\Rea\bigl(\chi(z_0)/\psi'(z_0)\bigr)>1/8$; (iv) $w(2)=(A^T+\tfrac12I)^{-1}\one>0$.
Then $\alpha=-1/2$, so the graph is never limit-stable at a legal $\beta$, yet its full support is a stable positive fixed point at
legal points with $\eps$ arbitrarily small (and $\beta\to2^+$).
\end{lemma}
\begin{proof}[Sketch]
By Theorem~\ref{thm:AS}, $\nu_A(c)=\tfrac12+R/c+O(c^{-2})$, while the legality boundary is
$c-\sqrt{c^2-c}=\tfrac12+\tfrac1{8c}+O(c^{-2})$. If $R>1/8$ there is room between them for all large $c$; choose a rational
$\gamma$ in between. Positivity holds near $\beta=2$ by (iv) and is $t$-free.
\end{proof}

Let $H=$ \dsix+&FNDL~[uds?+ (one of the $123$) and let $E_a$ be its blow-up with class sizes $(1,1,2,a,1,1,1)$, $n=7+a$; $E_0$ is
\dsix+&FNFL~[udo?+.

\begin{theorem}[Finite-$\eps$-only family; proved, refereed]\label{thm:famE}
For every $a\ge0$, $E_a$ is not a clique, $\alpha(\perpc{E_a})=-1/2$ exactly, and the full support of $E_a$ is a stable positive
fixed point at legal points with $\eps$ arbitrarily small. Thus for every $n\ge7$ there is a counterexample on $n$ nodes that
is invisible in the $\eps\to0$ limit.
\end{theorem}
\begin{proof}[Sketch]
Lemma~\ref{lem:W} gives $\Phi=z(z^2+z+1)\bigl((a+7)z(z^2+z+1)+2\bigr)$ identically in $a$ (non-edge picture); the cubic factor has all roots
in $\Rea z>-1/2$ iff $a+7>16/3$; in the adjacency picture, at $z_0=\tfrac12-\tfrac{\sqrt3}{2}i$ one has $1/\hat s\,'(z_0)=\tfrac14-\tfrac{\sqrt3}{12}i$, so $R^*=1/4>1/8$
for every $a$; and the class vector at $\beta=2$ is $(6,2,2,18,6,2,10)/(18a+17)>0$. Apply Lemma~\ref{lem:Bmech}.
Instances $a=0,\dots,5$ are certified at $t=20$, $\beta=821/400$, $(\eps,\delta)=(1/21,421/8400)$, $q=421/420$.\footnote{\cf{notes/research_wf3/theory/family_E_certificate_part1.json}, \cf{family_E_witnesses.json}.}
\end{proof}
The two-parameter engine applied to $E_1$ returns one stable cell, over the same $\beta$-cell and with the upper $t$-end at
$\beta=53/26$ in the same isolating bracket $(53,54)$ as for $E_0$.\footnote{\cf{certificates/n8_instances_region2d.json}.} The two
regions are not identical: at $\beta=53/26$, $t=107/2$ the full support of $E_0$ is stable and that of $E_1$ is not (Hurwitz
determinant~7 negative), so the upper end of the window depends on $a$. The cube-root factor arises from the root
locus of the class size: for the blow-up family of $H_3$ one has $\Phi_k(\omega)=(k-1)\omega^2$ at a primitive cube root of
unity $\omega$, so the locus passes through $\omega$ exactly at class size $2$. A float-steered scan of all one-node extensions of the
$124$ seven-node graphs, followed by exact triage of the $9{,}808$ extensions it flagged (an eigenvalue of $\perpc{A}$ near a
primitive cube root of unity), found $15$ eight-node graphs that satisfy every hypothesis of Lemma~\ref{lem:Bmech}, each with an
exact legal witness; extensions not flagged by the float prefilter were not examined exactly.\footnote{\cf{notes/research_wf3/theory/cube_n8_witnesses.json}.}

\subsection{Growing thresholds}
\begin{theorem}[Families Q, Q2, Y; proved, refereed]\label{thm:QQ2Y}
(Q) For the core \dsix+&F^S{^ReL|?+ with one sink class of size $k$ ($n=7+k$), every member is positive at every $\beta>2$, and
is limit-stable at $\beta=5/2$, $27/10$, $14/5$ for all $k\ge25$, $79$, $286$ respectively. So for every $n\ge32$ ($86$,
$293$) there are counterexamples on $n$ nodes with $q$ arbitrarily close to $3/2$ ($17/10$, $9/5$); along the family the limit
value of $q$ tends to $1.84416\ldots$.
(Q2) For the core \dsix+&F^U|m^QTq?+ with two sink classes of sizes $k$ and $3k$ ($n=7+4k$), limit polynomial
$z^3(2z^2+z+2)^2/4$, every member is positive at every $\beta>2$, and is limit-stable at $\beta=5/2$, $3$, $7/2$, $19/5$ for all
$k\ge1$, $3$, $10$, $33$.
(Y) For the eight-node core \dsix+&G^ZZZr|KjqTo+ with two sink classes of sizes $4k$ and $k$ ($n=8+5k$), every member is positive
at every $\beta>2$, and is limit-stable at $\beta=3$, $7/2$, $4$, $9/2$, $24/5$ for all $k\ge13$, $24$, $46$, $117$, $363$; the
family ceiling is $\beta\approx4.9607$.
\end{theorem}
Each threshold is proved sharp by exact real-root isolation in $k$ of the Hurwitz determinants of the shifted limit polynomial.
The eight-node core of (Y) is not among the hand-verified graphs.\footnote{\cf{notes/research_wf3/theory/family_Q_certificate.json}, \cf{family_Q2_certificate.json}, \cf{family_Y_certificate.json}.}
Certified instances include $n=32$ at $q=15000/10001$ (Q), $n=19$ at $q=20000/10001$ and $n=47$ at $q=25000/10001$ (Q2), and
$n=73$ at $q=20000/10001$ (Y). Exact lower bounds at small $n$ found by search and certified pointwise are listed in
Table~\ref{tab:qlow}.

\begin{table}[ht]
\centering
\small
\begin{tabular}{@{}rlll@{}}
\toprule
$n$ & graph & $(\eps,\delta)$ & $q$ (exact)\\
\midrule
8 & \dsix+&G^ZnRba\ba^w+ & $(1/1001,\ 1709/1001000)$ & $1709/1001\approx1.7073$\\
9 & \dsix+&H^\us{\HrEa^|~_+ & $(1/1001,\ 17691/10010000)$ & $17691/10010\approx1.7673$\\
10 & \dsix+&I^Qz|NbkfuNa^}~|tw+ & $(1/10001,\ 9209/50005000)$ & $18418/10001\approx1.8416$\\
11 & \dsix+&J^R\~R{\a^k^a^~^~]n^~?+ & $(1/100001,\ 861/45455000)$ & $17220/9091\approx1.8942$\\
\bottomrule
\end{tabular}
\caption{Certified lower bounds $q_n^*\ge q$ (not hand-verified).}\label{tab:qlow}
\end{table}

\subsection{\texorpdfstring{$q_n^*\to\infty$}{q*_n -> infinity}}
A fixed core with sink classes cannot push the limit threshold past the number of core nodes: the limit matrix
$D-\bar b\one^T$ has trace $\le-1$ when no class is a join. The way around this is a family of cores whose limit polynomial can be
prescribed.

\begin{theorem}[Unbounded thresholds; proved, refereed]\label{thm:U}
For every integer $N\ge3$ there is $n_0(N)$ such that for every $n\ge n_0(N)$ some non-complete $n$-node digraph has a stable
positive full-support fixed point at the legal points $(\eps,\delta)=(1/(1+t),(N-1)/(1+t))$ for all large $t$. Hence
$q_n^*\ge N-1$ for $n\ge n_0(N)$, and $q_n^*\to\infty$.
\end{theorem}
\begin{proof}[Sketch]
\emph{Arms forest.} In the non-edge digraph $D$, take a directed $3$-cycle $c_0\to c_1\to c_2\to c_0$ and attach ``arms'' $c_2\to a\to z_a\to(\text{target})$,
so that every node other than $c_2$ has exactly one $D$-out-neighbour and reaches $c_2$ (an in-tree rooted at $c_2$), while $c_2$
sends to $c_0$ and to the first node of each arm. By the matrix determinant lemma (Lemma~T of the
notes), $\chi_{D-e_{c_2}\one^T}(z)=z^{m-d}p(z)$ for any $p=z^d+z^{d-1}+c_2z^{d-2}+\dots+c_d$ with integers $c_2\ge1$, $c_j\ge0$
satisfying a mild forest constraint; the core has $m=2p(1)-3$ nodes. Lemma~Q (Sherman--Morrison) gives $\psi_D$ in closed form and
Lemma~P gives positivity in closed form. A single sink class receiving from every core node except $c_2$ has
$\Phi_k=z^{m-d}(G+kp)$ with $G$ fixed.
\emph{Target polynomial.} With $s=1/(2N)$ and an integer $W$, round the coefficients of
$\prod_{k=1}^N\bigl((z+s)^2+W^2k^2\bigr)$ (except the first two) to the nearest integer. By Rouch\'e (Lemma~S), for
$W>W_0(N)=2N\,4^N(N+1)^{2N-2}$ all roots of the rounded polynomial $\tilde p$ satisfy $-3/(4N)<\Rea z<-1/(4N)$, and all its
coefficients from $c_2$ on are $\ge1$.
\emph{Positivity} at $\beta=N$ holds for large $W$ because the top coefficient dominates.
\emph{Class size.} $\Phi_k/(m+k)\to\chi_L=z^{m-d}\tilde p$, so by root continuity there is $k_0$ with all roots of $\Phi_k$ in
$\Rea z>-1/N$ for $k\ge k_0$. Lemma~\ref{lem:limit} at the rational $\beta=N$ finishes, with $n_0=m+k_0$.\footnote{Theorem~U, \cf{notes/theorem_U/}; referee record \cf{notes/theorem_U_workflow.json}.}
\end{proof}
The bound $n_0(N)$ is super-exponential, and $k_0$ is not explicit in general. Exact reduced certificates give $n_0=1{,}812$,
$6{,}214{,}995$ and $1{,}590{,}755{,}528$ for $N=3,4,5$; a $65$-node core with $p$-coefficients $(1,1,11,7,10,3,1)$ gives
$q_n^*\ge3$, $4$, $9/2$, $23/5$ for all $n\ge219$, $657$, $2{,}475$, $4{,}909$. Two explicit graphs are certified directly by the
pointwise verifier without using any lemma: $n=28$ at $q=20000/10001$ and $n=41$ at $q=22500/10001$; at $k_0-1$ the explicit limit
test fails with a negative Hurwitz determinant, so the reduced thresholds are sharp there.\footnote{\cf{certificates/theorem_U_explicit_instances.json}.}
The only known upper bound remains $q_n^*\lesssim n-2$ from~\eqref{eq:geneson-bound}.

\subsection{Unbounded support size at one parameter point}
Theorem~\ref{thm:famF} uses a $k$-dependent schedule $\beta\to2$. At a fixed legal point the family $F_k$ eventually dies: at
$(\eps,\delta)=(1/21,11/210)$ it is stable exactly for $k=0,\dots,40$ and unstable for every $k\ge41$, and in fact at every fixed
legal point for large $k$, because $\chi_L$ has the factor $z^2+z+1$ with roots on $\Rea z=-1/2$. A different core avoids this.

Let $R=$ \dsix+&F^U|m^QTq?+ (one of the $124$; rows \texttt{0111110} \texttt{1011011} \texttt{1101101} \texttt{1100111} \texttt{1101001}
\texttt{0010101} \texttt{1100100}),
$S=(1,1,0,1,1,1,1)$ and $R_k:=R[S,\text{all};k]$, $n=7+k$.

\begin{theorem}[One point, every size; proved, refereed]\label{thm:A}
At $(\eps,\delta)=(1/21,11/210)$, i.e.\ $t=20$, $\beta=21/10$, $q=22/21$, the full support of $R_k$ is a stable positive non-clique
fixed point for every $k\ge0$. Consequently, at this point a non-complete $n$-node digraph with a stable positive full-support fixed
point exists if and only if $n\ge6$.
\end{theorem}
\begin{proof}[Sketch]
In the non-edge picture, (3) holds iff every eigenvalue of $D'+aJ$, $a=t/\beta$, has real part $>-1/\beta$, and for sink-class
families $\det(zI-aJ-D')=z^{p-1}\Psi(z)$ with $\Psi=z\det(zI-D-aJ)-a\sum_jk_j\det(zI-D+b_j\one^T)$. Here $\Psi_k=z^3g_k$ with
$g_k(z)=z^5-3z^2-2z-1-a[(k+7)z^4+(k+12)z^3+(3k+21)z^2+(k+13)z+(k+5)]$. At $a=200/21$, after the shift $z=-10/21-x$, every Hurwitz
determinant is a polynomial in $k$ with only positive coefficients (for example $21\Delta_1=200k+1450$), hence positive for all
$k\ge0$. Positivity: $w_R(21/10)>0$ and $\sigma_S=381654/295819>1$, so Lemma~\ref{lem:X}(X2) applies for every $k$. The
``only if'' part is Theorem~\ref{thm:n5}; $n=6$ is $G_1$.\footnote{Theorem~A, \cf{notes/theorem_A_fixed/}; \cf{certificates/theorem_A_fixed_parameter_instances.json} ($k=0,\dots,15$); referee record \cf{notes/theorem_A_fixed_workflow.json}.}
\end{proof}
An independent second proof (equitable-partition quotient, separate exact code) confirms the result; members up to $n=107$
were certified pointwise. Also proved in the notes (the referee spot-checked the first statement at $28$ exact legal points and
re-proved the self-join statement independently):
$R_k$ is an SFP at \emph{every} legal point with $\beta\le43/20$, for every $k\ge0$; every legal point with $\beta<4.9607$ carries an
infinite family (via Theorem~\ref{thm:QQ2Y}(Y)); and, conditionally on the lemmas of Theorem~\ref{thm:U}, every legal point
does, with an exact instance at $(\eps,\delta)=(1/2,2)$ for all $k\ge1581$. The $k$-fold self-join of $R$ ($n=7k$), which is not a
blow-up, is an SFP at every legal point with $\beta\le21/10$ for every $k\ge1$.

\begin{proposition}[A refuted conjecture; certified]\label{prop:P}
It is not true that positivity of $(D^T+rI)^{-1}\one$ for some $r\in(0,1/2)$ forces $D$ to have an eigenvalue with real part
$\le-1/2$. For example, the non-edge matrix of \dsix+&F^U|m^QTq?+ has every eigenvalue with real part $>-10/21$ while
$(D^T+\tfrac{10}{21}I)^{-1}\one>0$; and $33$ of the $54$ catalogue entries with $S$ equal to all nodes, e.g.\ the base \dsix+&FMAlW\d\y?+ with
$\chi_D=z^7-z^5-7z^4-10z^3-9z^2-4z-1$, have every eigenvalue of $D$ in $\Rea z>-1/2$ and a positive vector at $r=49/100$. In
particular the joins $G\vee K_k$ of such bases are limit-stable for every $k\ge1$.
\end{proposition}
(This was ``Conjecture P'' of the project notes; see Question~\ref{q:char}.)

%======================================================================
\section{Verification and reproducibility}\label{sec:verification}

\subsection{Certificates}
Every number in Sections~\ref{sec:n6}--\ref{sec:families} is read from a certificate file that can be regenerated by a named
script (Appendix~\ref{app:certs}). Values are stored exactly: integers, rationals, and algebraic numbers as an irreducible
integer minimal polynomial together with a rational isolating interval; decimals appear only as labelled approximations. A
region certificate stores, for each cell, the isolating data of its walls, an exact sample, the verdict with its reason (for
example ``Hurwitz determinant~6 negative''), and the outcome of the pointwise cross-check. The ``never'' records of the large
sweeps store critical values and cell counts rather than full decompositions, so verifying them means re-running the committed
engine; the records of all stable graphs are self-contained. Bulk files carry SHA-256 checksums.

\subsection{Independent audits}
The engine and the $n=6$ map were audited adversarially, including a from-scratch symbolic re-derivation of the region of $H_3$
(same set and same minimal polynomial). The $n=7$ prefilter and resolver were audited twice ($979{,}800$ exact comparisons; all
$786{,}596$ undecided lines recomputed independently, byte-identical), and the ray reduction, curve test and edge bound once (an
independent all-minors computation on $3{,}129$ graphs and $6{,}000$ random exact legal points on never-stable graphs). The family
theorems (Theorems~\ref{thm:famF}, \ref{thm:famE}, \ref{thm:QQ2Y}, \ref{thm:U}, \ref{thm:A}), the $626$-entry catalogue and the
threshold tool were each refereed by a separate reviewer who re-derived the steps, taking Lemma~\ref{lem:limit}(b),(c) as given; no blocking
issue was found. Non-fatal defects found
by the referees (e.g.\ a misprinted node count, $2p(1)-1$ for $2p(1)-3$, in the notes on Theorem~\ref{thm:U}) are logged and
corrected here.

\subsection{Human verification}
The two new six-node counterexamples were hand-verified on 14 October 2026, and all $124$ seven-node counterexamples on
15 October 2026; the project owner recorded the latter hand check as not fully confident and made sign-off conditional on
the machine re-verification below. The $124$ were additionally re-decided by four exact verifiers written from the definitions without access
to the project's code: one with exact rationals and its own Hurwitz determinants; one with pure fractions and its own
elimination and characteristic polynomials (with the transpose convention reversed, $0$ of the $124$ survive, a useful
convention check); one by exact right-half-plane root counting without Hurwitz determinants (Routh array, Sturm sequences and
exact isolation, required to agree); and one checking structural claims (pairwise non-isomorphism, domination-freeness,
non-completeness, the exact limit test and $\alpha(E_0)=-1/2$). All $124$ were confirmed at $742$ legal points with $0$ undecided
cases.\footnote{\cf{notes/verify_124_workflow.json}; \cf{certificates/n7_candidates_legal_witnesses.json}.} The verification also
found that $15$ of the first stored witnesses in the limit file are illegal points ($q\le1$); a separate file lists one legal
witness per graph. Graphs with $n\ge8$ nodes have not been hand-verified.

\subsection{Nondegeneracy}
The conjecture is usually stated for nondegenerate CTLNs. At a witness point of each of the three six-node and $124$ seven-node
counterexamples we certified exactly that every one of the $2^n-1$ supports has $\det(I-W_\sigma)\ne0$, a fixed point with no zero
coordinate and no off-support drive equal to $0$, and that the full support is stable there.\footnote{\cf{certificates/nondegeneracy_witnesses.json}, script \cf{n7/nondegeneracy_check.py}.}

\subsection{Machine-checked statements}
Lean~4 (toolchain v4.35.0) with Mathlib (tag v4.35.0) checks $87$ theorems with no \texttt{sorry}; the axiom audit shows only
\texttt{propext}, \texttt{Classical.choice} and \texttt{Quot.sound}. They comprise: the closed form~\eqref{eq:closed} with the
transpose and $I-W=\eps K$; the legality equivalence~\eqref{eq:legal}; Hurwitz stability of $-I+W$ iff positive stability of $K$;
the ray identity, characteristic-polynomial scaling and propagation of stability down rays (Lemma~\ref{lem:rayF}); the
intermediate-value step with the explicit root, and the forward direction of Lemma~\ref{lem:rayR}; Lemma~\ref{lem:A2};
Lemma~\ref{lem:e2} including its spectral step; Lemma~\ref{lem:D} with its induced-subgraph form; Stodola's condition
(Lemma~\ref{lem:stodola}); and the twin eigenvalue (Lemma~\ref{lem:twin}). Not formalised: the converse of Lemma~\ref{lem:rayR}
(it needs continuity of eigenvalues), the limit criterion Lemma~\ref{lem:limit}, the rank-one identities of
Lemma~\ref{lem:rankone}(ii)--(v), the off-support condition, and the full certificate lemma ``(1)--(3) imply a stable equilibrium of
\eqref{eq:tln} with support $\sigma$''.\footnote{\cf{notes/lean/README.md}, \cf{notes/lean/CtlnCert/}.}

\paragraph{Trust base.} The certificates rely on: exact polynomial arithmetic of a standard computer-algebra system for the projection polynomials of the
decomposition (determinants, factorisation, gcd, resultants, discriminants; the projection identities were re-verified by
independent fraction-free determinants on a sample of $321$ records); the Routh--Hurwitz theorem, which is cited; nauty's
\texttt{geng}/\texttt{directg} for the seven-node enumeration, validated by its count; and the C prefilters, validated by
reproducing the six-node data exactly.

%======================================================================
\section{Open problems}\label{sec:open}

\begin{question}[Growth rate]\label{q:growth}
How fast does $q_n^*$ grow? Theorem~\ref{thm:U} gives $q_n^*\to\infty$ with a super-exponential $n_0(N)$; the certified points
are $q$ close to $2$ at $n=19$, $5/2$ at $n=47$ and $23/5$ for $n\ge4{,}909$, while the only upper bound is Geneson's $n-2$. Is
$q_n^*\ge c\log n$, or $q_n^*\ge n^c$? The bottleneck is integer polynomials with $c_1=1$, nonnegative coefficients, roots in a thin
strip $\Rea z>-1/N$, and small value at $1$.
\end{question}

\begin{question}[The $n=8$ map]
Complete the $\eps\to0$ classification at $n=8$. An overflow-guarded spectral prefilter reproduces all $n\le7$ data exactly; the
full pass is estimated at $500$--$800$ CPU-hours, and besides the family members of Section~\ref{sec:families} one further eight-node
limit-stable graph, \dsix+&GLjlNZ}NzH\o+ ($41$ edges; certified, not hand-verified), was found in a random dense sample. The finite-$\eps$ map at $n=8$ ($1.79$ trillion digraphs) is out of reach by enumeration alone.
\end{question}

\begin{question}[Formalisation]
Formalise the converse of Lemma~\ref{lem:rayR} (openness of positive stability), the limit criterion Lemma~\ref{lem:limit}, and the
certificate lemma (1)--(3)$\Rightarrow$ stable equilibrium, so that every certified case becomes a checked instance.
\end{question}

\begin{question}[Characterising limit-stable graphs]\label{q:char}
Proposition~\ref{prop:P} refutes the natural spectral obstruction. Which graphs are limit-stable, i.e.\ have $w(\gamma)>0$ and
$\alpha<-\gamma$ at a common $\beta>2$? A related question: is $\alpha\le-1/2$ for every graph whose full support is an SFP at some
legal point? This holds for $n\le7$ (Corollary~\ref{cor:alpha7}), with equality attained for every $n\ge7$ by Theorem~\ref{thm:famE};
no proof is known.
\end{question}

\begin{question}[Proper supports]
The truth of the conjecture depends only on full supports (Theorem~\ref{thm:reduction}), but listing all pairs $(G,\sigma)$
realising a counterexample needs condition~(2): describe the admissible in-sets for the $n\le7$ counterexamples. At $n=7$ all
$12{,}288$ pairs formed by a six-node counterexample and a seventh node with any in/out pattern have been certified
($7{,}744$ stable somewhere, $4{,}544$ never, $0$ escalations; \cf{certificates/n7_proper_supports_all.jsonl}; not
hand-verified); a structural description is open.
\end{question}

%======================================================================
\appendix
\section{Certificates}\label{app:certs}
Paths are relative to the project archive. ``Script'' is the program that regenerates the file.

\begingroup\footnotesize
\begin{longtable}{@{}p{0.18\textwidth}p{0.49\textwidth}p{0.26\textwidth}@{}}
\toprule
claim & certificate & script\\
\midrule
\endhead
Geneson's graph reproduced & \cf{certificates/stage0_geneson_2607.21396.json} & test suite, block 8\\
$n=5$ package, diff zero & \cf{certificates/stage1_n5_diff.json} & \cf{stage1_n5_diff.py}\\
engine validation & \cf{certificates/validate_fastregion_n4.json}, \cf{_n5.json}, \cf{region2d_v2_validation.json} & \cf{validate_fastregion.py}, \cf{region2d.py}\\
$k\le4$ & \cf{certificates/region2d_n2to4_all.json} & \cf{region2d.py}\\
Theorem~\ref{thm:n5} & \cf{certificates/region2d_n5_shard0of1*}, \cf{n5_summary.json} & \cf{region2d.py}\\
Theorem~\ref{thm:n6} & \cf{certificates/n6_summary.json}, \cf{region_n6_limit_*}, \cf{region2d_n6_shard*}, \cf{region_n6_eps*}, \cf{COUNTEREXAMPLE_CANDIDATE_n6_idx*.json} & \cf{fastregion.py}, \cf{region2d.py}\\
twin classes & \cf{certificates/n6_survivor_equivalence_classes.json} & \cf{notes/prover_scripts_J/}\\
Theorem~\ref{thm:q6} & \cf{certificates/qstar/qstar_n6.json} & \cf{n7/qstar.py}\\
Table~\ref{tab:chain} & \cf{certificates/n7_finite_summary.json}, \cf{n7_curve_TOTAL.json}, \cf{n7_region2d/} & \cf{n7/run_n7_curve.py}, \cf{run_n7_region2d.py}, \cf{aggregate_n7.py}\\
limit map $n=7$ & \cf{certificates/COUNTEREXAMPLE_CANDIDATES_n7_limit.jsonl}, \cf{n7_limit_tally.json} & \cf{n7/run_n7_limit.py}\\
the $124$ graphs & \cf{certificates/COUNTEREXAMPLE_CANDIDATES_n7_finite.jsonl}, \cf{n7_candidates_legal_witnesses.json} & \cf{n7/aggregate_n7.py}, \cf{notes/verify_124/}\\
$E_0$ & \cf{certificates/COUNTEREXAMPLE_CANDIDATE_n7_finite_FNFL.json}, \cf{qstar/qstar_fnfl.json} & \cf{n7/write_finite_certificates.py}, \cf{n7/qstar.py}\\
Proposition~\ref{prop:Mfalse} & \cf{notes/lemma_M_refuted.md}, \S2 & \cf{notes/prover_scripts_M/}\\
Proposition~\ref{prop:bulge} & \cf{certificates/n7_bulge_Rstar_exact.json}, \cf{n7_limit_hits_finite_eps_bulge_witnesses.json} & \cf{notes/bulge_exact/}\\
Theorem~\ref{thm:q7} & \cf{certificates/qstar/qstar_n7_finite.json}, \cf{qstar_n7limit.json} & \cf{n7/qstar.py --sharpen}\\
Theorems~\ref{thm:famF}, \ref{thm:famE}, \ref{thm:QQ2Y}, Table~\ref{tab:qlow} & \cf{notes/research_wf3/theory/family_*_certificate.json}, \cf{family_E_witnesses.json}, \cf{top_n8.json}, \cf{cert_q_910.json}, \cf{cert_q_11.json}, \cf{families_catalogue_n{6,7}.json} & \cf{family_*_proof.py}, \cf{certify_top.py}, \cf{certify_q.py}, \cf{run_family_prover.py}, \cf{verify_all.py}\\
eight-node instances & \cf{certificates/n8_instances_region2d.json}, \cf{notes/research_wf3/theory/cube_n8_witnesses.json} & \cf{n7/n8_instances_region2d.py}, \cf{witness_cube.py}\\
Theorem~\ref{thm:U} & \cf{certificates/theorem_U_explicit_instances.json}, \cf{notes/theorem_U/} & \cf{theorem_U.py}, \cf{certify_forest.py}, \cf{certify_explicit.py}\\
Theorem~\ref{thm:A} & \cf{certificates/theorem_A_fixed_parameter_instances.json}, \cf{notes/theorem_A_fixed/} & \cf{theorem_A.py}, \cf{independent_A.py}\\
nondegeneracy & \cf{certificates/nondegeneracy_witnesses.json} & \cf{n7/nondegeneracy_check.py}\\
Lean & \cf{notes/lean/CtlnCert/} & \texttt{lake build}, \texttt{Audit.lean}\\
\bottomrule
\end{longtable}
\endgroup

{\small
\bibliographystyle{plain}
\bibliography{refs}
}

\end{document}
